\beginchapter Chapter 16. Typing\\Math Formulas

\TeX\ is designed to handle complex ^{mathematical expressions} in such a way
that most of them are easy to input. The basic idea is that a complicated
formula is composed of less complicated ^{formulas} put together in a
simple way; the less complicated formulas are, in turn, made up of simple
combinations of formulas that are even less complicated; and so on. Stating
this another way, if you know how to type simple formulas and how to combine
formulas into larger ones, you will be able to handle virtually any formula
at all. So let's start with simple ones and work our way up.

The simplest formula is a single letter, like `$x$', or a single number,
like `2'. In order to put these into a \TeX\ text, you type `|$x$|' and
`|$2$|', respectively. Notice that all mathematical formulas are enclosed
in special math brackets; we are using |$| as the math bracket in this
manual, in accord with the plain \TeX\ format defined in Appendix~B\null,
because mathematics is supposedly expensive.

When you type `|$x$|' the `$x$' comes out in italics, but when you
type `|$2$|' the `$2$' comes out in roman type.  In general, all
characters on your keyboard have a special interpretation in math
formulas, according to the normal conventions of mathematics printing:
Letters now denote ^{italic} letters, while digits and punctuation
denote ^{roman} digits and punctuation; a hyphen ({\tt-}) now denotes a ^{minus
sign} ($-$), which is almost the same as an em-dash but not quite (see
Chapter~2). The first |$| that you type puts you into ``^{math mode}'' and the
second takes you out (see Chapter~13).  So if you forget one |$| or type
one |$| too many, \TeX\ will probably become thoroughly confused and you
will probably get some sort of error message. ^^{dollarsign}

Formulas that have been typeset by a printer who is unaccustomed to
mathematics usually look quite strange to a mathematician, because a
novice printer usually gets the spacing all wrong. In order to alleviate
this problem, \TeX\ does most of its own spacing in math formulas; and it
{\sl ignores\/} any ^{spaces} that you yourself put between |$|'s. For
example, if you type `|$ x$|' and `|$ 2 $|', they will mean the same thing
as `|$x$|' and `|$2$|'. You can type `\hbox{|$(x + y)/(x - y)$|}' or
`|$(x+y) / (x-y)$|', but both will result in `$(x+y)/(x-y)$', a formula in
which there is a bit of extra space surrounding the $+$ and~$-$ signs but
none around the~/~sign.  Thus, you do not have to memorize the complicated
rules of math spacing, and you are free to use blank spaces in any way you
like.  Of course, spaces are still used in the normal way to mark the end
of control sequences, as explained in Chapter~3.  In most circumstances
\TeX's spacing will be what a mathematician is accustomed to; but we will
see in Chapter~18 that there are control sequences by which you can
override \TeX's spacing rules if you want to.

One of the things mathematicians like to do is make their formulas look
like ^{Greek} to the uninitiated. In plain \TeX\ language you can type
`|$$\alpha, \beta, \gamma, \delta;$$|' and you will get the first four
Greek letters ^^|\alpha|^^|\beta|^^|\gamma|^^|\delta|
$$\alpha,\beta,\gamma,\delta;$$ furthermore there are uppercase Greek
letters like `$\Gamma$', which you can get by typing `|$\Gamma$|'.
^^|\Gamma| Don't feel intimidated if you aren't already familiar with Greek
letters; they will be easy to learn if you need them. The only difficulty
is that some symbols that look nearly the same must be carefully
distinguished. For example, the Greek letters ^|\nu|~($\nu$) and
^|\kappa|~($\kappa$) should not be confused with the italic letters $v$
and~$x$; the Greek ^|\phi|~($\phi$) is different from the slashed zero
called ^|\emptyset|~($\emptyset$).  A~lowercase epsilon ($\epsilon$) is
quite different from the symbol used to denote membership in a set
($\in$); type `|$\epsilon$|' for $\epsilon$ and `|$\in$|' for $\in$.
^^|\epsilon| ^^|\in| Some of the lowercase Greek letters have variant
forms in plain \TeX's math italic fonts:
`|$(\phi,\theta,\epsilon,\rho)$|' yields
`$(\phi,\theta,\epsilon,\rho)$' while
`|$(\varphi,\vartheta,\varepsilon,\varrho)$|' yields
`$(\varphi,\vartheta,\varepsilon,\varrho)$'.
^^|\phi|^^|\theta|^^|\rho|^^|\varphi|^^|\vartheta|^^|\varrho|^^|\varepsilon|

Besides Greek letters, there are a lot of ^{funny symbols} like `$\approx$'
(which you get by typing `|$\approx$|') ^^|\approx|^^{special symbols for math}
and `$\mapsto$' (which you ^^|\mapsto|^^{math symbols}
get by typing `|$\mapsto$|').
 A complete list of these control sequences
and the characters they correspond to appears in Appendix~F\null.
Such control sequences are allowed only in math mode, i.e., between
|$|'s, because the corresponding symbols appear in the math fonts.

\exercise What should you type to get the formula
`$\gamma+\nu\in\Gamma$'\thinspace?
\answer |$\gamma+\nu\in\Gamma$|.

\exercise Look at Appendix F to discover the control sequences for
`$\le$', `$\ge$', and~`$\ne$'. \ (These are probably the three most
commonly used math symbols that are not present on your keyboard.)
\ What does plain \TeX\ call them?
\answer ^|\le|, ^|\ge|, and ^|\ne|. \ (These are short for ``less-or-equal,''
``greater-or-equal,'' and ``not-equal.'') \ You can also use the names
^|\leq|, ^|\geq|, and ^|\neq|. \ (The fourth most common symbol is, perhaps,
`$\infty$', which stands for ``^{infinity}'' and is called `^|\infty|'.)

Now let's see how the more complex formulas get built up from simple ones.
In the first place, you can get ^{superscripts} $\rm^{(up\,high)}$ and
^{subscripts} $\rm_{(down\,low)}$
^^{indices, see subscripts}
^^{superiors, see superscripts} ^^{inferiors, see subscripts}
by using `|^|' and `|_|', as shown in the following examples:
\beginmathdemo
\it Input&\it Output\cr
\noalign{\vskip2pt}
|$x^2$|&x^2\cr
|$x_2$|&x_2\cr
|$2^x$|&2^x\cr
|$x^2y^2$|&x^2y^2\cr
|$x ^ 2y ^ 2$|&x ^ 2y ^ 2\cr
|$x_2y_2$|&x_2y_2\cr
|$_2F_3$|&_2F_3\cr
\endmathdemo
Notice that |^| and |_| apply only to the next single character.
If you want several things to be superscripted or subscripted, just enclose
them in braces:
\beginmathdemo
|$x^{2y}$|&x^{2y}\cr
|$2^{2^x}$|&2^{2^x}\cr
|$2^{2^{2^x}}$|&2^{2^{2^x}}\cr
|$y_{x_2}$|&y_{x_2}\cr
|$y_{x^2}$|&y_{x^2}\cr
\endmathdemo
The braces in these examples have been used to specify ``^{subformulas},''
i.e., simpler parts of a larger formula.  \TeX\ makes a box for each
subformula, and treats that box as if it were a single symbol.  Braces
also serve their usual purpose of grouping, as discussed in Chapter~5.

It is illegal to type `|x^y^z|' or `|x_y_z|'; \TeX\ will
complain of a ``double superscript'' or ``double subscript.'' You must type
`|x^{y^z}|' or `|x^{yz}|' or
`|x_{y_z}|' or `|x_{yz}|' in order to make your intention clear.

A superscript or subscript following a character applies to that character
only; but when following a subformula it applies to that whole subformula,
and it will be raised or lowered accordingly. For example,
\beginmathdemo
|$((x^2)^3)^4$|&((x^2)^3)^4\cr
|${({(x^2)}^3)}^4$|&{({(x^2)}^3)}^4\cr
\endmathdemo
In the first formula the `|^3|' and `|^4|' are superscripts on the ^{right
parentheses}, i.e., on the `|)|' characters that immediately precede them,
but in the second formula they are superscripts on the subformulas that
are enclosed in braces. The first alternative is preferable, because it is
much easier to type and it is just as easy to read.

\danger A subscript or superscript following nothing (as in the `|_2F_3|'
example on the preceding page, where the `|_2|' follows nothing) is taken
to mean a subscript or superscript of an empty subformula. Such notations
are (fortunately) rare in mathematics; but if you do encounter them it is
better to make your intention clear by showing the empty subformula
explicitly with braces.  In other words, the best way to get `${}_2F_3$'
in a formula is to type `|{}_2F_3|' or `|{_2}F_3|' or `|{_2F_3}|'.

\dangerexercise What difference, if any, is there between the output of
`|$x + _2F_3$|' and the output of `|$x + {}_2F_3$|'\thinspace?
\answer In the former, the `|_2|' applies to the plus sign ($x + _2F_3$);
but in the latter, it applies to an empty subformula ($x + {}_2F_3$).

\dangerexercise Describe the differences between the outputs of `|${x^y}^z$|'
and `|$x^{y^z}$|'.
\answer The results are `${x^y}^z$' and `$x^{y^z}$'; the $z$ in the first
alternative is the same size as the $y$, but in the second it is smaller.
Furthermore, the $y$ and $z$ in the first case aren't quite at the same height.
\ (Good typists never even think of the first construction, because
mathematicians never want it.)

You can have simultaneous subscripts and superscripts, and you can specify them
in any order:
\beginmathdemo
|$x^2_3$|&x^2_3\cr
|$x_3^2$|&x_3^2\cr
|$x^{31415}_{92}+\pi$|&x^{31415}_{92}+\pi\cr
\noalign{\smallskip}
|$x_{y^a_b}^{z_c^d}$|&x_{y^a_b}^{z_c^d}\cr
\endmathdemo
Notice that simultaneous su$\rm_b^{per\kern-1pt}$scripts are positioned
over each other.  However, a subscript will be ``tucked in'' slightly when it
follows certain letters; for example, `|$P_2^2$|' produces `$P_2^2$'.
If for some reason you want the left edges of both subscript and superscript
to be aligned, you can fool \TeX\ by inserting a null subformula:
`|$P{}_2^2$|' produces `$P{}_2^2$'.

The control sequence ^|\prime| stands for the symbol `$\prime$', which
is used mostly in superscripts. In fact, `$\prime$' is so big as it stands
that you would never want to use it except in a subscript or superscript,
where it occurs in a smaller size. Here are some typical examples:
\beginmathdemo
\it Input&\it Output\cr
\noalign{\vskip2pt}
|$y_1^\prime$|&y_1^\prime\cr
|$y_2^{\prime\prime}$|&y_2^{\prime\prime}\cr
|$y_3^{\prime\prime\prime}$|&y_3^{\prime\prime\prime}\cr
\endmathdemo
Since single and double primes occur rather frequently, plain \TeX\
provides a convenient abbreviation: You can simply type |'| instead
of |^\prime|, and |''| instead of |^{\prime\prime}|, and so on.
\beginmathdemo
|$f'[g(x)]g'(x)$|&f'[g(x)]g'(x)\cr
|$y_1'+y_2''$|&y_1'+y_2''\cr
|$y'_1+y''_2$|&y'_1+y''_2\cr
|$y'''_3+g'^2$|&y'''_3+g'{}^2\cr
\endmathdemo

\dangerexercise Why do you think \TeX\ treats |\prime| as a large symbol
that appears only in superscripts, instead of making it a smaller
symbol that has already been shifted up into the superscript position?
\answer The second alternative doesn't work properly when there's a
subscript at the same time as a prime. Furthermore, some mathematicians
use |\prime| also in the subscript position; they write, for example,
$F'(w,z)=\partial F(w,z)/\partial z$ and $F_\prime(w,z)=\partial F(w,z)/
\partial w$.

\dangerexercise Mathematicians sometimes use ``^{tensor notation}''
in which subscripts and superscripts are staggered, as in `$R_i{}^{jk}{}_l$'.
Explain how to achieve such an effect.
\answer |$R_i{}^{jk}{}_l$|.

Another way to get complex formulas from simple ones is to use the control
sequences ^|\sqrt|, ^|\underline|, or ^|\overline|.
^^{surds, see sqrt} ^^{vinculum, see overline}
Like |^| and |_|, these operations apply to the character or subformula
that follows them:
\beginmathdemo
|$\sqrt2$|&\sqrt2\cr
|$\sqrt{x+2}$|&\sqrt{x+2}\cr
|$\underline4$|&\underline4\cr
|$\overline{x+y}$|&\overline{x+y}\cr
|$\overline x+\overline y$|&\overline x+\overline y\cr
|$x^{\underline n}$|&x^{\underline n}\cr
|$x^{\overline{m+n}}$|&x^{\overline{m+n}}\cr
|$\sqrt{x^3+\sqrt\alpha}$|&\sqrt{x^3+\sqrt\alpha}\cr
\endmathdemo
You can also get cube roots `$\root3\of{\phantom{h}}$' and similar things
by using ^|\root|:
\beginmathdemo
|$\root 3 \of 2$|&\root 3 \of 2\cr
|$\root n \of {x^n+y^n}$|&\root n \of {x^n+y^n}\cr
|$\root n+1 \of a$|&\root n+1 \of a\cr
\endmathdemo

\danger The |\sqrt| and |\underline| and |\overline| operations are able to
place lines above or below subformulas of any size or shape; the bar lines
change their size and position, so that they are long enough to cover the
subformula, and high enough or low enough not to bump into it. For example,
consider `|\overline|~|l|' ($\,\overline l\,$) versus `|\overline|~|m|'
($\,\overline m\,$):  The first has a shorter bar line, and this line has
been raised higher than the bar in the second. Similarly, the bar in
`|\underline|~|y|' ($\,\underline y\,$) is lower than the bar in
`|\underline|~|x|' ($\,\underline x\,$); and square root signs appear in a
variety of positions based on the height and depth of what is being
|\sqrt|'d: $\sqrt a + \sqrt d + \sqrt y$. \TeX\ knows the height, depth, and
width of every letter and every subformula, because it considers them to be
boxes, as explained in Chapter~11. If you have a formula in which there is
only one |\sqrt|, or only one |\overline| or |\underline|, the normal
positioning rules work fine; but sometimes you want to have uniformity
between different members of a complex formula. For example, you might want to
typeset `$\mskip-1mu\sqrt{\mathstrut a}+
 \sqrt{\mathstrut d}+\sqrt{\mathstrut y}\mskip+2mu$',
putting all square roots in the same vertical position. There's an easy way
to do this, using the control sequence ^|\mathstrut| as follows:
\begintt
$\sqrt{\mathstrut a}+\sqrt{\mathstrut d}+\sqrt{\mathstrut y}$.
\endtt
A |\mathstrut| is an invisible box whose width is zero; its height and depth
are the height and depth of a parenthesis `('. Therefore subformulas
that contain |\mathstrut| will always have the same height and depth,
unless they involve more complicated constructions like subscripts and
superscripts. Chapter~18 discusses more powerful operations called ^|\smash|
and ^|\phantom| by which you can obtain complete control over the positioning
of roots and similar signs.

\exercise Test your understanding of what you have read so far in this chapter
by explaining what should be typed to get the following formulas. \ (Be
sure to check your answer with Appendix~A to confirm that you're right.)
$$\hbox to\hsize{\indent
  $\displaystyle
  10^{10}\hfil 2^{n+1}\hfil (n+1)^2\hfil \sqrt{1-x^2}\hfil
  \overline{w+\overline z}\hfil p_1^{e_1}\hfil a_{b_{c_{d_e}}}\hfil
  \root3\of{h''_n(\alpha x)}\hfil$}$$
\answer |10^{10}|; \stretch|2^{n+1}|; \stretch|(n+1)^2|; \stretch|\sqrt{1-x^2}|;
\stretch|\overline{w+\overline z}|; \stretch|p_1^{e_1}|; \stretch
|a_{b_{c_{d_e}}}|; \stretch|\root3\of{h''_n(\alpha x)}|.
\ (Of course, you should enclose these formulas in dollar signs so that
\TeX\ will process them in math mode.  Superscripts and subscripts can be
given in either order; for example, |h''_n| and |h_n''| both work the
same.  You should not leave out any of the braces shown here; for example,
`|$10^10$|' would yield `$10^10$'. But it doesn't hurt to insert
additional braces around letters or numbers, as in `|({n}+{1})^{2}|'. The
indicated blank spaces are necessary unless you use extra braces;
otherwise \TeX\ will complain about undefined control sequences
|\overlinez| and |\alphax|.)

\exercise What mistake did B. C. ^{Dull} discover after he typed the
following?
\begintt
If$ x = y$, then $x$ is equal to $y.$
\endtt
\answer He got `If$ x = y\ldots$' because he forgot to leave a space
after `|If|'; ^{spaces} disappear between dollar signs. He should
also have ended the sentence with `|$y$.|'; punctuation that belongs
to a sentence should not be included in a formula, as we will see
in Chapter~18. \ (But you aren't expected to know that yet.)

\exercise Explain how to type the following sentence:
\begindisplay
Deleting an element from an $n$-tuple leaves an $(n-1)$-tuple.
\enddisplay
\answer |Deleting an element from an $n$-tuple leaves an $(n-1)$-tuple.|

\exercise List all the italic letters that descend below the baseline.
\ (These are the letters for which |\underline| will lower its bar line.)
\answer $Q,f,g,j,p,q,y$. \ (The analogous ^{Greek} letters are
^^{italic letters with descenders} ^^{descenders}
$\beta,\gamma,\zeta,\eta,\mu,\xi,\rho,\phi,\varphi,\chi,\psi$.)

We have discussed the fact that the characters you type have special meanings
in math mode, but the examples so far are incomplete; they don't reveal all
the power that is at your fingertips just after you press the `|$|' key.
It's time now to go back to basics: Let us make a systematic survey of
what each character does, when it is used in a formula.

The 52 ^{letters} (|A| to |Z| and |a| to |z|) denote italic symbols
($A$~to~$Z$ and $a$~to~$z$), which a mathematician would call ``^{variables}.''
\TeX\ just calls them ``^{ordinary symbols},'' because they make up the
bulk of math formulas. There are two variants of lowercase L in plain \TeX,
namely `$l$' (which you get by simply typing `|l|') and `$\ell$'
(which you get by typing `^|\ell|'). Although mathematicians commonly
write something that looks like `$\ell$' in their manuscripts, they
do so only to distinguish it from the numeral~`1'.  This
distinguishability problem is not present in printed mathematics, since an
italic `$l$' is quite different from a~`1'; therefore it is traditional to
use `$l$' unless `$\ell$' has been specifically requested.

Plain \TeX\ also treats the 18 characters
\begintt
0 1 2 3 4 5 6 7 8 9 ! ? . || / ` @ "
\endtt
^^{digits}^^{numerals}
as ordinary symbols; i.e., it doesn't insert any extra space when these
symbols occur next to each other or next to letters. Unlike the letters,
these 18 characters remain in roman type when they appear in formulas.
There's nothing special for you to remember about them, except that the
^{vertical line} `\|' has special uses that we shall discuss later.
Furthermore, you should be careful to distinguish between `oh' and `zero':
The italic letter~$O$ is almost never used in formulas unless it appears
just before a left parenthesis, as in `$O(n)$'; ^^{big-$O$ notation}
and the numeral~$0$ is almost never used just before a left parenthesis
unless it is preceded by another digit, as in `$10(n-1)$'. Watch for
left parentheses and you'll be $0K$. \ (Lowercase o's also tend to
appear only before left parentheses; type `|x_0|' instead of `|x_o|',
since the formula `$x_0$' is generally more correct than `$x_o$'.)

The three characters |+|, |-|, and |*|
are called ``^{binary operations},'' because they operate on two parts of
a formula. For example, |+|~is a ^{plus sign}, which is used for the sum
of two numbers; |-|~is a ^{minus sign}. The ^{asterisk}~(|*|) is rarer
^^{star, see asterisk}
in mathematics, but it also behaves as a binary operation. Here are
some examples of how \TeX\ typesets binary operations when they appear
next to ordinary symbols:
\beginmathdemo
\it Input&\it Output\cr
\noalign{\vskip2pt}
|$x+y-z$|&x+y-z\cr
|$x+y*z$|&x+y*z\cr
|$x*y/z$|&x*y/z\cr
\endmathdemo
Notice that |-| and |*| produce quite different math symbols from what you
get in normal text: The ^{hyphen}~(-) becomes a minus sign~($-$), and
the raised asterisk~(*) drops down to a lower level~($*$).

\danger \TeX\ does not treat |/| as a binary operation, even though a
^{slash} stands for division (which qualifies as a binary operation on
mathematical grounds). The reason is that printers traditionally put extra
space around the symbols $+$, $-$, and~$*$, but not around~$/$. If \TeX\ were to
typeset |/| as a binary operation, the formula `|$1/2$|' would come out
`$1\mathbin/2$', which is wrong; so \TeX\ considers |/| to be an ordinary
symbol.

\danger Appendix F lists many more binary operations, for which you type
control sequences instead of single characters. Here are some examples:
\beginmathdemo
|$x\times y\cdot z$|&x\times y\cdot z\cr
|$x\circ y\bullet z$|&x\circ y\bullet z\cr
|$x\cup y\cap z$|&x\cup y\cap z\cr
|$x\sqcup y\sqcap z$|&x\sqcup y\sqcap z\cr
|$x\vee y\wedge z$|&x\vee y\wedge z\cr
|$x\pm y\mp z$|&x\pm y\mp z\cr
\endmathdemo
It is important to distinguish $\times$ (^|\times|) from $X$ (|X|)
and from $x$ (|x|); to distinguish $\cup$ (^|\cup|) from $U$ (|U|)
and from $u$ (|u|); to distinguish $\vee$ (^|\vee|) from $V$ (|V|)
and from $v$ (|v|); to distinguish $\circ$ (^|\circ|) from $O$ (|O|)
and from $o$ (|o|). ^^|\cdot|^^|\bullet|^^|\cap|^^|\sqcup|^^|\sqcap|^^|\wedge|
^^{cross, see dagger, times}
^^|\pm|^^|\mp| The symbols `$\lor$' and `$\land$' can also be called
^|\lor| and ^|\land|, since they frequently stand for binary operations
that are called ``^{logical~or}'' and ``^{logical~and}.''

\danger Incidentally, binary operations are treated as ordinary symbols
if they don't occur between two quantities that they can operate on.
For example, no extra space is inserted next to the $+$, $-$, and~$*$
in cases like the following:
\beginmathdemo
|$x=+1$|&x=+1\cr
|$3.142-$|&3.142-\cr
|$(D*)$|&(D*)\cr
\endmathdemo
Consider also the following examples, which show that binary
operations can be used as ordinary symbols in superscripts and subscripts:
\beginmathdemo
|$K_n^+,K_n^-$|&K_n^+,K_n^-\cr
|$z^*_{ij}$|&z^*_{ij}\cr
|$g^\circ \mapsto g^\bullet$|&g^\circ \mapsto g^\bullet\cr
|$f^*(x) \cap f_*(y)$|&f^*(x) \cap f_*(y)\cr
\endmathdemo

\dangerexercise How would you obtain the formulas `$z^{*2}$'
and `$h_*'(z)$'\thinspace?
\answer |$z^{*2}$| and |$h_*'(z)$|.

Plain \TeX\ treats the four characters |=|, |<|, |>|, and |:|\ as
``^{relations}'' because they express a relationship between two
quantities. For example, `${x<y}$' means that $x$~is less than~$y$.
Such relationships have a rather different meaning from binary
operations like $+$, and the symbols are typeset somewhat differently:
\beginmathdemo
|$x=y>z$|&x=y>z\cr
|$x:=y$|&x:=y\cr
|$x\le y\ne z$|&x\le y\ne z\cr
|$x\sim y\simeq z$|&x\sim y\simeq z\cr
|$x\equiv y\not\equiv z$|&x\equiv y\not\equiv z\cr
|$x\subset y\subseteq z$|&x\subset y\subseteq z\cr
\endmathdemo
^^|\le|^^|\ne|^^|\simeq|^^{colon}^^{equals}^^{lessthan}^^{greaterthan}
^^{colonequals}^^|\equiv|^^|\not|^^|\subset|^^|\subseteq|^^|\sim|
^^{hooks, see subset, supset} ^^{wiggle, see sim}
(The last several examples show some of the many other
relational symbols that plain \TeX\ makes available via control sequences;
see Appendix~F.)

The two characters `|,|' (^{comma}) and `|;|' (^{semicolon}) are treated
as ^{punctuation marks in formulas}; this means that \TeX\ puts a little
extra space after them, but not before them.
\beginmathdemo
|$f(x,y;z)$|&f(x,y;z)\cr
\endmathdemo
It isn't customary to put extra space after a `|.|'\ (^{period}) in
math formulas, so \TeX\ treats a period as an ordinary symbol.
If you want the `|:|'\ character to be treated as a punctuation mark
instead of as a relation, just call it ^|\colon|:
\beginmathdemo
|$f:A\to B$|&f:A\to B\cr
|$f\colon A\to B$|&f\colon A\to B\cr
\endmathdemo
^^|\to|
If you want to use a comma as an ordinary symbol (e.g., when it
appears in a large number), just put it in braces; \TeX\ treats
anything in braces as an ordinary symbol. For instance,
\beginmathdemo
|$12,345x$|&12,345x\qquad\rm(wrong)\cr
|$12{,}345x$|&12{,}345x\qquad\,\rm(right)\cr
\endmathdemo

\dangerexercise What's an easy way to get a raised dot in a decimal
constant (e.g., `$3{\cdot}1416$')?
\answer |$3{\cdot}1416$|. \ (One of the earlier examples in this
chapter showed that ^|\cdot| is a binary operation; putting it in braces
makes it act like an ordinary symbol.)\par
If you have lots of constants like this, for example in a table, there's a way
to make ordinary periods act like |\cdot| symbols: Just define
^|\mathcode||`.| to be |"0201|, assuming that the fonts of plain \TeX\ are
being used.  However, this could be dangerous, since ordinary
periods are used frequently in displayed equations; the |\mathcode| change
should be confined to places where every period is to be a |\cdot|.

So far we have considered letters, other ordinary symbols, binary operations,
relations, and punctuation marks; hence we have covered almost every key on
the typewriter. There are just a few more: The characters `|(|' ^^{lparen}
and `|[|' ^^{lbracket} are called ``^{openings},'' while `|)|' ^^{rparen} and
^^{fences, see opening, closing, delimiters}
`|]|' ^^{rbracket} are called ``^{closings}''; these act pretty much like
ordinary symbols, but they help \TeX\ to decide when a binary operation is
not really being used in a binary way.  Then there is the character~|'|,
which we know is used as an abbreviation for |\prime| superscripts.
Finally, we know that plain \TeX\ reserves the other ten characters:
\begintt
\ $ % # & ~ { } _ ^
\endtt
These are not usable for symbols in math mode unless their ^|\catcode|
values are changed (see Chapter~7). Although |{| and |}| specify
grouping, the control sequences `|\{|' and `|\}|' ^^|\lbrace|^^|\rbrace| can
be used to get `$\{$' as an opening and `$\}$' as a closing.

\ddanger All of these math mode interpretations are easily changeable, since
each character has a ^|\mathcode|, as explained in Chapter~17; none of
the conventions are permanently built into \TeX\null. However, most of them are
so standard that it is usually unwise to make many changes, except perhaps
in the interpretations of |`|, |"|, and |@|. ^^{at sign} ^^{leftquote}
^^{doublequote}

The special characters |^| and |_| that designate superscripts
^^{circumflex}^^{underbar}
and subscripts should not be used except in formulas. Similarly,
the names of math symbols like |\alpha| and |\approx|, and the
control sequences for math operations like |\overline|, must not
invade ordinary text. \TeX\ uses these facts to detect ^{missing dollar
signs} in your input, before such mistakes cause too much trouble. For
example, suppose you were to type
\begintt
The smallest $n such that $2^n>1000$ is~10.
\endtt
\TeX\ doesn't know that you forgot a `|$|' after the first `|n|', because
it doesn't understand English; so it finds a ``formula'' between the
first two |$| signs:
\begindisplay
The smallest $n such that $
\enddisplay
after which it thinks that `|2|' is part of the text. But then the |^|
reveals an inconsistency; \TeX\ will automatically insert a~|$| before
the~|^|, and you will get an error message. In this way the computer has
gotten back into synch, and the rest of the document can be typeset as if
nothing had happened.

\danger Conversely, a blank line or ^|\par| is not permitted in math mode.
This gives \TeX\ another way to recover from a missing~|$|; such
errors will be confined to the paragraph in which they occur.

\danger If for some reason you cannot use |^| and |_| for superscripts
and subscripts, because you have an unusual keyboard or because you need
|^| for French accents or something, plain \TeX\ lets you type ^|\sp|
and ^|\sb| instead. For example, `|$x\sp2$|' is another way to get `$x\sp2$'.
On the other hand, some people are lucky enough to have keyboards that
contain additional symbols besides those of standard ASCII.
^^{character set}
When such symbols are available, \TeX\ can be set up to make math
typing a bit more pleasant. For example, at the author's installation there
are keys labeled \up\ and~\dn\ that produce visible symbols
(these make superscripts and subscripts look much nicer
on the screen); there are keys for the relations {\tentex\char'34},
{\tentex\char'35}, and {\tentex\char'32} (these save time); and there are
about two dozen more keys that occasionally come in handy. \ (See Appendix~C.)
^^{uparrow}^^{downarrow}^^{leq}^^{geq}^^{neq}

\danger Mathematicians are fond of using ^{accents} over letters, because
this is often an effective way to indicate relationships between
mathematical objects, and because it greatly extends the number of available
symbols without increasing the number of necessary fonts.
Chapter~9 discusses the use of accents in ordinary text, but mathematical
accents are somewhat different, because spacing is not the same; \TeX\ uses
special conventions for accents in formulas, so that the two sorts of
accents will not be confused with each other. The following math accents
are provided by plain~\TeX:
\beginmathdemo
|$\hat a$|&\hat a\cr
|$\check a$|&\check a\cr
|$\tilde a$|&\tilde a\cr
|$\acute a$|&\acute a\cr
|$\grave a$|&\grave a\cr
|$\dot a$|&\dot a\cr
|$\ddot a$|&\ddot a\cr
|$\breve a$|&\breve a\cr
|$\bar a$|&\bar a\cr
|$\vec a$|&\vec a\cr
\endmathdemo
^^|\hat|^^|\check|^^|\tilde|^^|\acute|^^|\grave|
^^|\dot|^^|\ddot|^^|\breve|^^|\bar|^^|\vec|
The first nine of these are called |\^|, |\v|, |\~|, |\'|, |\`|, |\.|,
|\"|, |\u|, and |\=|, respectively, when they appear in text; |\vec| is an
accent that appears only in formulas. \TeX\ will complain if you try to
use |\^| or |\v|, etc., in formulas, or if you try to use |\hat| or
|\check|, etc., in ordinary text.

\danger It's usually a good idea to define special control sequences for
accented letters that you need frequently. For example, you can put
^^|\def|
\begintt
\def\Ahat{{\hat A}}
\def\chat{{\hat c}}
\def\scheck{{\check s}}
\def\xtilde{{\tilde x}}
\def\zbar{{\bar z}}
\endtt
at the beginning of a manuscript that uses the symbols $\hat A$, $\hat c$,
$\check s$, $\tilde x$, and $\bar z$ more than, say, five times.  This
saves you a lot of keystrokes, and it makes the manuscript easier to
read.  Chapter~20 explains how to define control sequences.

\danger When the letters $i$ and $j$ are accented in math formulas, ^{dotless}
symbols $\imath$ and $\jmath$ should be used under the accents. These
symbols are called ^|\imath| and ^|\jmath| in plain \TeX. Thus, for example,
a paper that uses `$\hat\imath$' and `$\hat\jmath$' ought to begin with
the following definitions:
\begintt
\def\ihat{{\hat\imath}}
\def\jhat{{\hat\jmath}}
\endtt

\def\Ahat{{\hat A}}
\danger You can put ^{accents on top of accents}, making symbols like
$\skew6\hat\Ahat$ that might cause a mathematician to
squeal with ecstasy. However, it takes a bit of finesse to get the
upper accent into a position that looks right, because the designer of a
font for mathematics usually tells \TeX\ to position math accents in
special ways for special letters. Plain \TeX\ provides a control sequence
called ^|\skew| that makes it fairly easy to shift superaccents into their
proper place. For example, `|\skew6\hat\Ahat|' was used to produce the
symbol above. The number `|6|' in this example was chosen by trial and
error; `|5|'~seems to put the upper accent a bit too far left, while
`|7|'~makes it a bit too far right, at least in the author's opinion.
The idea is to fiddle with the amount of skew until you find what
pleases you best.

\danger It's possible, in fact, to put math accents on any subformula, not
just on single characters or accented characters. But there's usually not
much point in doing so, because \TeX\ just centers the accent over the
whole subformula. For example, `|$\hat{I+M}$|' yields `$\hat{I+M}$'. In
particular, a |\bar| accent always stays the same size; it's not like
^|\overline|, which grows with the formula under it. Some people prefer
the longer line from |\overline| even when it applies to only a single
letter; for example, `|$\bar z+\overline z$|' produces `$\bar z+\overline
z$', and you can take your pick when you define |\zbar|.  However, plain
\TeX\ does provide two accents that grow; they are called ^|\widehat| and
^|\widetilde|:
\beginmathdemo
|$\widehat x,\widetilde x$|&\tenmath\widehat x,\widetilde x\cr
|$\widehat{xy},\widetilde{xy}$|&\tenmath\widehat{xy},\widetilde{xy}\cr
|$\widehat{xyz},\widetilde{xyz}$|&\tenmath\widehat{xyz},\widetilde{xyz}\cr
\endmathdemo
The third example here shows the maximum size available.

\def\ghat{{\hat g}}
\exercise This has been another long chapter; but cheer up, you have learned
a lot! Prove it by explaining what to type in order to get the formulas
$e^{-x^2}$, $D\sim p^\alpha M+l$, and $\ghat\in(H^{\pi_1^{-1}})'$. \ (In
the last example, assume that a control sequence |\ghat| has already been
defined, so that |\ghat| produces the accented letter $\ghat$.)
\answer |$e^{-x^2}$|, |$D\sim p^\alpha M+l$|, and |$\ghat\in(H^{\pi_1^{-1}})'$|.
\ (If you are reading the dangerous bend sections, you know that the
recommended way to define |\ghat| is `|\def\ghat{{\hat g}}|'.)

\endchapter

Producing ^{Greek} letters is as easy as $\pi$.
You just type |... as easy as $\pi$.|
\author LESLIE ^{LAMPORT}, {\sl The ^{L\kern-.2em\raise.6ex\hbox{a}%
  \kern-.1em\TeX} Document Preparation System\/} (1983)
        % Note: the final manual has a slightly different wording on p43.
        % It's now called "LaTeX: A Document Preparation System" (1986)
        % But I decided to cite the original, partly because I have
        % no smallcaps sans-serif `A' to match the new LaTeX logo!

\bigskip

\TeX\ has no regard for the glories of the Greek tongue---\/
as far as it is concerned, Greek letters are just additional weird symbols,
and they are allowed\/ {\rm only} in math mode.
In a pinch you can get the output $\tau\epsilon\chi$ by typing %
  |$\tau\epsilon\chi$|,
but if you're actually setting Greek text, you will be using
a different version of \TeX, designed for a keyboard with Greek letters on it,
and you shouldn't even be reading this manual,
which is undoubtedly all English to you.
\author MICHAEL ^{SPIVAK}, {\sl The Joy of \TeX\/} (1982)

\eject
